% TEMPLATE-MILC-v3.tex - plantilla maestra UNICA de las guias MILC v3.
%
% La usan las cuatro familias de guias, cada una con su driver:
%   · Grados 6-11          -> scripts/build-guias-g11.py
%   · Bebras               -> scripts/build-guias-bebras.py
%   · Semillero            -> scripts/build-guias-semillero.py
%   · Territorio interior  -> scripts/build-guias-territorio-interior.py
%
% Antes habia tres copias casi identicas (~400 lineas iguales, 6 distintas):
% cada mejora habia que aplicarla tres veces, y en la practica no se hacia
% (el motor de recursos vivia solo en dos de ellas). Lo que varia entre
% familias esta parametrizado en 6 placeholders que cada driver llena
% (se nombran aqui SIN su sintaxis real: el builder reemplaza por texto plano,
% tambien dentro de los comentarios, y un valor multilinea rompe el preambulo):
%
%   HEADER_IZQ        encabezado izquierdo de pagina
%   HEADER_DER        encabezado derecho de pagina
%   PORTADA_ETIQUETA  rotulo sobre el numero grande (GRADO/MOMENTO/MODULO)
%   PORTADA_SERIE     linea de serie de la portada
%   PORTADA_ID        identificador de la guia en la portada
%   PORTADA_CIRCULO   el circulito de grado (°), o vacio si no aplica
%   PDF_TITULO        titulo en texto plano (sin \textbf ni comillas TeX) para
%                     los metadatos del PDF (/Title); lo produce lib_xelatex.texto_pdf
%
% Composicion (auditoria 2026-09-03): las bandas de estacion piden espacio
% (\Needspace*) y prohiben el salto justo despues (\nopagebreak), asi no quedan
% huerfanas al pie; las cajas largas (softbox, drawbox, infoband, puentebox) son
% `breakable`, asi una caja de media pagina ya no arrastra todo a la siguiente
% dejando la anterior al 40 %. Todo texto sobre mostaza o verde va en milcNegro
% (blanco sobre #E5B400 da 1,93:1 y sobre #58A923 2,95:1; ninguno pasa WCAG AA).
% El resto de claves las documenta content/guias/_SCHEMA.md. El script valida
% que no queden placeholders sin reemplazar antes de escribir el archivo final.

% Etiquetado PDF (latex-lab): /Lang, estructura y texto alternativo de las
% figuras, para lectores de pantalla. Debe ir ANTES de \documentclass.
\DocumentMetadata{lang=es-CO,tagging=on}
\documentclass[11pt,letterpaper]{article}
% headheight=24pt: la cabecera derecha lleva dos lineas (identidad de la guia
% y estacion actual). headsep se acorta en la misma medida para que el bloque
% de texto quede exactamente donde estaba con headheight=12pt.
\usepackage[margin=2.0cm,headheight=24pt,headsep=13pt]{geometry}
\usepackage[table]{xcolor}
\usepackage{fontspec}
\usepackage[spanish]{babel}
\usepackage{tikz}
% Librerías TikZ para los diagramas de las guías (bloque `recursos:`):
%  · babel  → babel en español vuelve activos caracteres como «>», y TikZ los usa
%             en sus opciones (>=stealth, ->). Sin esta librería, cualquier
%             diagrama con flechas revienta con «\language@active@arg> has an
%             extra }». Debe ir ANTES de que se dibuje nada.
%  · positioning        → sintaxis «below left=2pt and 3pt».
%  · decorations.pathreplacing → llaves ({brace}) para señalar rangos.
\usetikzlibrary{babel,positioning,decorations.pathreplacing}
\usepackage{graphicx}
% Circuitos: el trabajo de electrónica se hace hoy con micro:bit y sus sensores
% integrados (MakeCode, sin cableado), así que no se necesitan esquemáticos.
% Si algún día entran montajes con protoboard, basta descomentar esta línea:
% los diagramas .tex/.tikz declarados en `recursos:` ya se incrustan solos.
% \usepackage{circuitikz}
\usepackage[most]{tcolorbox}
\usepackage{tabularx}
\usepackage{array}
\usepackage{fancyhdr}
\usepackage{titlesec}
\usepackage[document]{ragged2e}  % bandera a la derecha en todo el cuerpo
\usepackage{enumitem}
\usepackage{microtype}
\usepackage{etoolbox}
\usepackage{needspace}
% hyperref al final del preambulo: metadatos del PDF (titulo, autor, idioma,
% familia), marcadores por estacion (\pdfbookmark) y enlaces sin recuadro.
\usepackage[hidelinks,
  pdftitle={El acuerdo de aula: la norma que uno se da a sí mismo},
  pdfauthor={PhD. Álvaro Cárdenas Orozco},
  pdfsubject={Territorio interior · Educación socioemocional · Grado 9° · Momento 10 · MILC v3},
  pdflang={es-CO},
  bookmarksopen=true,bookmarksnumbered=false]{hyperref}

% Helvetica Neue no trae la flecha U+2192, asi que cada «→» escrito literalmente
% en el YAML se compilaba a NADA, en silencio (462 flechas en 61 guias: rutas del
% tipo «paleta → tipografia → lockup» salian rotas). Mapeamos el caracter a su
% version matematica, que si tiene glifo. Asi el autor puede seguir escribiendo
% «→» comodamente en el YAML.
\usepackage{newunicodechar}
\newunicodechar{→}{\ensuremath{\rightarrow}}
% Mismo problema con los simbolos de marca y vinetas: el «✓» de «una palomita ✓»
% salia como cuadrito vacio en el PDF de todas las guias que lo usaban.
\usepackage{pifont}
\newunicodechar{✓}{\ding{51}}
\newunicodechar{✗}{\ding{55}}
\newunicodechar{✔}{\ding{52}}
\newunicodechar{•}{\textbullet}
\newunicodechar{≥}{\ensuremath{\geq}}
\newunicodechar{≤}{\ensuremath{\leq}}

\setmainfont{Helvetica Neue}
\setsansfont{Helvetica Neue}
\setlength{\parindent}{0pt}
\setlength{\parskip}{6pt}
% Sin particion de palabras: en 6.º-7.º una palabra cortada por guion al final
% de linea («Comuni-cacion») frena la lectura mucho mas que una linea corta.
\hyphenpenalty=10000
\exhyphenpenalty=10000
\pretolerance=2000
\tolerance=3000
\setlength{\emergencystretch}{3em}
\linespread{1.10}
\renewcommand{\arraystretch}{1.25}
\setlist{leftmargin=5mm,itemsep=1mm,topsep=1mm}
\newcolumntype{Y}{>{\RaggedRight\arraybackslash}X}

\definecolor{milcMagenta}{HTML}{E6007E}
\definecolor{milcVino}{HTML}{5A0038}
\definecolor{milcTurquesa}{HTML}{0093A5}
\definecolor{milcVerde}{HTML}{58A923}
\definecolor{milcMostaza}{HTML}{E5B400}
\definecolor{milcOcre}{HTML}{B86600}
\definecolor{milcNegro}{HTML}{241816}
\definecolor{milcGris}{HTML}{F7F7F4}
\definecolor{milcLinea}{HTML}{E8E2DD}
\definecolor{milcGradePrimary}{HTML}{0D9488}
\definecolor{milcGradeDark}{HTML}{115E59}
\definecolor{milcGradeSoft}{HTML}{CCFBF1}
\definecolor{milcGradeAccent}{HTML}{CFE7FF}
\definecolor{milcGradeText}{HTML}{FFFFFF}
\color{milcNegro}

\pagestyle{fancy}
\fancyhf{}
% Cabecera de dos lineas: identidad de la guia arriba y, a la derecha y en
% gris, la estacion en curso (\markright desde \milcestacion). Antes de la
% primera estacion la segunda linea va vacia; el \strut mantiene la altura
% para que la linea de arriba no baile entre paginas.
\lhead{\small Territorio interior · Educación socioemocional\\[-1pt]{\footnotesize\strut}}
\rhead{\small Grado 9° · Momento 10 · MILC v3\\[-1pt]{\footnotesize\strut\color{milcNegro!65}\rightmark}}
\cfoot{\small\thepage}
\renewcommand{\headrulewidth}{0pt}

% Sobre mostaza (#E5B400) y verde (#58A923) el texto va en milcNegro; sobre el
% resto de la paleta, en blanco. Se usa dentro de fonttitle/fontupper, donde un
% \color posterior manda sobre coltitle.
\newcommand{\milctintasobre}[1]{%
  \ifboolexpr{test{\ifstrequal{#1}{milcMostaza}} or test{\ifstrequal{#1}{milcVerde}}}%
    {\color{milcNegro}}{\color{white}}}

\newtcolorbox{titlebox}[1]{
  enhanced,colback=#1,colframe=#1,arc=7mm,boxrule=0pt,
  left=7mm,right=7mm,top=3mm,bottom=3mm,
  before skip=8pt,after skip=8pt,
  fontupper=\sffamily\bfseries\Large\color{white}
}

\newtcolorbox{softbox}[3]{
  enhanced,breakable,pad at break=3mm,
  colback=#2,colframe=#1,borderline west={4pt}{0pt}{#1},
  arc=2mm,boxrule=.4pt,left=4mm,right=4mm,top=3mm,bottom=3mm,
  before skip=6pt,after skip=8pt,title={#3},coltitle=white,
  colbacktitle=#1,fonttitle=\sffamily\bfseries\milctintasobre{#1},fontupper=\RaggedRight,
  attach boxed title to top left={xshift=3mm,yshift=-2mm},
  boxed title style={arc=2mm, boxrule=0pt}
}

\newtcolorbox{drawbox}[2]{
  enhanced,breakable,pad at break=4mm,
  colback=white,colframe=#1,arc=4mm,boxrule=1pt,
  left=5mm,right=5mm,top=4mm,bottom=4mm,before skip=8pt,after skip=8pt,
  title={#2},coltitle=white,colbacktitle=#1,fonttitle=\sffamily\bfseries\milctintasobre{#1},
  fontupper=\RaggedRight,
  attach boxed title to top left={xshift=4mm,yshift=-2mm},
  boxed title style={arc=3mm, boxrule=0pt}
}

\newtcolorbox{infoband}[1]{
  enhanced,breakable,pad at break=2mm,colback=#1!8,colframe=#1!50,arc=2mm,boxrule=.5pt,
  left=4mm,right=4mm,top=2mm,bottom=2mm,before skip=4pt,after skip=4pt,
  fontupper=\small\RaggedRight
}

% Puente narrativo entre secciones (contrato editorial Conéctate).
% Da continuidad: "Acabas de X. Ahora vas a Y porque Z."
% Visualmente: banda izquierda en color del grado, texto en itálica pequeña.
\newtcolorbox{puentebox}{
  enhanced,breakable,pad at break=2mm,colback=milcGradeSoft,colframe=milcGradeDark,
  borderline west={3pt}{0pt}{milcGradeDark},
  arc=0mm,boxrule=0pt,left=5mm,right=4mm,top=2.5mm,bottom=2.5mm,
  before skip=4pt,after skip=8pt,
  fontupper=\itshape\small\color{milcNegro}\RaggedRight
}
% La flecha va en modo matematico a proposito: Helvetica Neue NO tiene el glifo
% U+2192, asi que \textrightarrow se compilaba a NADA («Missing character: There
% is no → in font Helvetica Neue») y el puente salia sin flecha. En modo
% matematico la flecha viene de la fuente matematica y si se dibuja.
\newcommand{\puente}[1]{%
  \begin{puentebox}%
    \textcolor{milcGradeDark}{$\rightarrow$}\hspace{2mm}#1%
  \end{puentebox}%
}

\newcommand{\linead}[1]{\noindent\color{milcLinea}\rule{#1}{0.5pt}\color{milcNegro}}
\newcommand{\respuesta}[1]{\par\vspace{2mm}\foreach \n in {1,...,#1}{\noindent\color{milcLinea}\rule{\linewidth}{0.5pt}\par\vspace{5mm}}\color{milcNegro}}
\newcommand{\checkbox}{\textcolor{milcGradeDark}{\fbox{\phantom{x}}}\hspace{5pt}}

% ─── Listas aireadas para las guías socioemocionales ────────────────────
% Componer las verificaciones como «\checkbox texto\\» deja el texto que salta
% de línea debajo del cuadrito y apelmaza el bloque. Estos entornos alinean la
% continuación y airean los ítems. Los usa Territorio Interior; cualquier
% familia puede hacerlo.
\newlist{checklista}{itemize}{1}
\setlist[checklista]{
  label=\textcolor{milcGradeDark}{\fbox{\phantom{x}}},
  leftmargin=9mm, labelsep=3.2mm, itemsep=2.4mm,
  topsep=2.5mm, parsep=0pt, partopsep=0pt, before=\RaggedRight
}
\newlist{pasoslista}{enumerate}{1}
\setlist[pasoslista]{
  label=\textcolor{milcGradeDark}{\sffamily\bfseries\arabic*.},
  leftmargin=8mm, labelsep=3mm, itemsep=2.8mm,
  topsep=1mm, parsep=0pt, partopsep=0pt, before=\RaggedRight
}

\newcommand{\phasecard}[4]{
\begin{tcolorbox}[enhanced,colback=#3,colframe=#2,arc=3mm,boxrule=.5pt,height=3.05cm,valign=center,halign=center,left=1.5mm,right=1.5mm,top=2mm,bottom=2mm]
{\sffamily\bfseries\fontsize{22}{24}\selectfont\textcolor{#2}{#1}}\par
{\sffamily\bfseries\small #4}
\end{tcolorbox}
}

% ── Piezas del rediseno 2026 (legibilidad 6.º-11.º) ───────────────────
% Banda de estacion: reemplaza el titlebox macizo. Titulo a la izquierda,
% dato practico (tiempo/modalidad) a la derecha, en una sola linea.
% Nunca huerfana: pide 9 lineas libres (o salta de pagina rellenando con
% \vfill) y prohibe el salto justo despues de la banda. Nueve y no seis:
% la caja partible que sigue necesita ~80pt (titulo adosado, relleno y dos
% lineas) para arrancar en la misma pagina; con menos, tcolorbox la manda
% entera a la siguiente y la banda se queda sola. El penalty 10000 va
% en `after`, ANTES del espacio posterior: un `after skip` (aunque sea 0pt)
% es un \vskip tras la caja, y ese glue es un punto de corte legal que un
% \nopagebreak escrito despues de \end{tcolorbox} ya no cierra. Cada banda
% deja marcador PDF y actualiza la cabecera derecha.
\newcounter{milcestacion}
\newcommand{\milcestacion}[2]{%
\Needspace*{9\baselineskip}%
\stepcounter{milcestacion}%
\pdfbookmark[1]{#2}{milcestacion\themilcestacion}%
\markright{#2}%
\begin{tcolorbox}[enhanced,colback=#1,colframe=#1,arc=2mm,boxrule=0pt,
  left=5mm,right=5mm,top=2.6mm,bottom=2.6mm,before skip=11pt,
  after={\par\nopagebreak\vskip7pt}]
{\sffamily\bfseries\fontsize{15}{18}\selectfont\milctintasobre{#1}#2}
\end{tcolorbox}}

% Tarjeta de paso numerada: sustituye las tablas «Pilar / Aplicacion», que
% eran formato de consulta docente, no de estudiante. El numero va en un
% circulo sobre el borde para que la secuencia se lea de un vistazo.
\newcommand{\milcpaso}[3]{%
\Needspace{4\baselineskip}%
\begin{tcolorbox}[enhanced,colback=white,colframe=milcLinea,arc=1.5mm,boxrule=.9pt,
  left=14mm,right=4mm,top=3mm,bottom=3mm,before skip=4pt,after skip=4pt,
  fontupper=\RaggedRight,
  overlay={\node[anchor=north west,circle,fill=#2,text=white,inner sep=0pt,
    minimum size=7.4mm,font=\sffamily\bfseries\small]
    at ([xshift=3.6mm,yshift=-3.2mm]frame.north west) {#1};}]
#3
\end{tcolorbox}}

% Casilla marcable de tamano util con lapiz (la anterior era \fbox{\phantom{x}},
% de alto variable segun el texto que la seguia).
\newcommand{\milcmarca}{\framebox[3.8mm]{\rule{0pt}{3.4mm}}}

% Fila marcable de lista suelta (checks de la estacion 1).
\newcommand{\milccheckrow}[1]{%
\par\vspace{1.8mm}\noindent
\makebox[6.4mm][l]{\raisebox{-.55mm}{\textcolor{milcNegro}{\milcmarca}}}%
\parbox[t]{\dimexpr\linewidth-6.4mm\relax}{\RaggedRight #1}\par}

% Celda marcable dentro de la rubrica (tabla de autoevaluacion). Misma
% sangria francesa, pero pensada para vivir dentro de una columna X.
\newcommand{\milcopcion}[1]{%
\makebox[5.2mm][l]{\raisebox{-.55mm}{\textcolor{milcNegro}{\milcmarca}}}%
\parbox[t]{\dimexpr\linewidth-5.2mm\relax}{\RaggedRight #1}}

% ── Recursos visuales de la guía ──────────────────────────────────────
% Declarados en el bloque `recursos:` del YAML y emitidos por el builder.
% \guiaFigura   → imagen rasterizada o PDF (foto, carta celeste, montaje).
% \guiaDiagrama → fuente TikZ/circuitikz (.tex) incrustada como vector:
%                 es la vía para circuitos y esquemáticos editables.
% [#1] = texto alternativo (alt) · #2 = ruta relativa al .tex · #3 = pie
% El alt llega al PDF etiquetado (/Alt) y lo lee el lector de pantalla.
\newcommand{\guiaFigura}[3][]{%
\begin{center}
\includegraphics[width=0.88\linewidth,height=0.38\textheight,keepaspectratio,alt={#1}]{#2}\\[1.5mm]
{\footnotesize\itshape\textcolor{milcNegro!75}{#3}}
\end{center}
}

\newcommand{\guiaDiagrama}[2]{%
\begin{center}
\input{#1}\\[1.5mm]
{\footnotesize\itshape\textcolor{milcNegro!75}{#2}}
\end{center}
}

\begin{document}
\thispagestyle{empty}

% ── Portada ───────────────────────────────────────────────────────────
% Antes: un rectangulo de color a pagina completa. Se veia bien en pantalla,
% pero en la fotocopiadora del colegio se comia el toner y salia gris sucio.
% Ahora la banda ocupa el tercio superior y el resto va en blanco.
\begin{tikzpicture}[remember picture,overlay]
  \path[fill=milcGradePrimary]
    (current page.north west) rectangle ([yshift=-10.6cm]current page.north east);

  \node[anchor=north west,text=milcGradeText] at ([xshift=2.0cm,yshift=-1.55cm]current page.north west)
    {\sffamily\bfseries\fontsize{12}{14}\selectfont MOMENTO};
  \node[anchor=north west,text=milcGradeText,opacity=.85] at ([xshift=2.0cm,yshift=-2.35cm]current page.north west)
    {\sffamily\bfseries\fontsize{8.5}{10}\selectfont TERRITORIO INTERIOR · SOCIOEMOCIONAL};
  \node[anchor=north east,text=milcGradeText,opacity=.85,align=right] at ([xshift=-2.0cm,yshift=-1.55cm]current page.north east)
    {\sffamily\bfseries\fontsize{9}{12}\selectfont I.E. Sor María Juliana\\Cartago, Valle del Cauca};

  \node[anchor=west,text=milcGradeText] (gradonum) at ([xshift=1.85cm,yshift=-7.1cm]current page.north west)
    {\sffamily\bfseries\fontsize{112}{112}\selectfont 10};
  % El circulito de grado (°) solo aplica a las guias de grado («11°»); en
  % Bebras y Semillero el numero grande es un momento/modulo, no un grado.
  % Se posiciona relativo al nodo `gradonum`, no en coordenadas absolutas.
  

  \node[anchor=west,text=milcGradeText,text width=11.4cm,align=left] at ([xshift=1.5cm,yshift=0.5cm]gradonum.east)
    {
     {\sffamily\bfseries\fontsize{9}{11}\selectfont\MakeUppercase{Territorio interior · Socioemocional}}\\[3mm]
     {\sffamily\bfseries\fontsize{21}{25}\selectfont\setlength{\baselineskip}{27pt}El acuerdo\\[4pt]de aula\\[4pt]la norma que uno se da}\\[3.5mm]
     {\sffamily\bfseries\fontsize{15}{18}\selectfont Grado 9° · Momento 10}
    };
\end{tikzpicture}

\vspace*{9.4cm}
\vfill

{\sffamily\bfseries\fontsize{8}{10}\selectfont\textcolor{milcGradeDark}{LO QUE TE LLEVAS HOY}}

\begin{tcolorbox}[enhanced,colback=white,colframe=milcNegro,arc=2mm,boxrule=1.6pt,
  left=5mm,right=5mm,top=4mm,bottom=4mm,before skip=3pt,after skip=12pt,
  fontupper=\RaggedRight\fontsize{11}{15.5}\selectfont]
El acuerdo de aula del curso: entre cinco y ocho reglas escritas por ustedes a partir de los nueve momentos del año, cada una con lo que pasa si se rompe, quién lo verifica y cuándo se revisa —firmado por todos y con fecha de revisión puesta—.
\end{tcolorbox}

\begin{tcolorbox}[enhanced,colback=milcGris,colframe=milcGris,arc=2mm,boxrule=0pt,
  left=5mm,right=5mm,top=3.5mm,bottom=3.5mm,after skip=12pt]
{\sffamily\bfseries\fontsize{8}{10}\selectfont\textcolor{milcNegro!55}{ESTA GUÍA ES DE}}\\[3mm]
\begin{tabularx}{\linewidth}{X p{3.2cm} p{3.2cm}}
\linead{\linewidth} & \linead{3.0cm} & \linead{3.0cm}\\[-1mm]
{\fontsize{8.5}{10}\selectfont\textcolor{milcNegro!55}{Nombre completo}} &
{\fontsize{8.5}{10}\selectfont\textcolor{milcNegro!55}{Grupo}} &
{\fontsize{8.5}{10}\selectfont\textcolor{milcNegro!55}{Fecha}}\\
\end{tabularx}
\end{tcolorbox}

\vfill

\noindent\textcolor{milcLinea}{\rule{\linewidth}{1pt}}\\[2mm]
{\sffamily\bfseries\fontsize{9.5}{12}\selectfont Autor: Dr. Álvaro Cárdenas Orozco}\\[1mm]
{\sffamily\fontsize{8.5}{11}\selectfont\textcolor{milcNegro!55}{Metodología MILC · Apertura ancestral · Norma propia y acuerdo · Triángulo Dussel-Estoicismo-Floridi}}

\newpage

\section*{Apertura: El acuerdo de aula: la norma que uno se da a sí mismo}

\begin{softbox}{milcGradeDark}{milcGradeSoft}{Saber ancestral}
En el Cauca, los pueblos organizados en el \textbf{CRIC} y el \textbf{pueblo misak} hacen algo que en un colegio suena casi imposible: \textbf{se dan sus propias normas}. \emph{No las reciben hechas: las deliberan en asamblea y las declaran como \textbf{mandatos}} ---la palabra ya lo dice todo, porque un mandato es lo que \textbf{la comunidad le manda} a sus autoridades y a sí misma---. \textbf{Y hay dos rasgos que nos importan hoy.} \textbf{Primero: quien hace la norma es quien va a vivir bajo ella.} \emph{No la escribe alguien de afuera que no va a sufrir sus consecuencias.} \textbf{Segundo: hay rendición de cuentas, y es pública.} \emph{Las autoridades ---que llevan el bastón que trabajaste en el momento 9 del grado 8--- \textbf{responden ante la asamblea que las nombró}, y lo hacen delante de todos.} \emph{Fíjate en el par completo, porque es lo que hace que el sistema funcione:} \textbf{la comunidad se da la norma \emph{y} vigila a quien la aplica}. \emph{Un curso que escribe reglas y no decide quién verifica ni ante quién se responde acaba de hacer la mitad del trabajo ---y es la mitad que no sirve sola---.}
\end{softbox}

\begin{softbox}{milcTurquesa}{milcGris}{Saber contemporáneo}
¿Y esto funciona, o es solo bonito? \textbf{Hay una investigación que le dio a su autora el Nobel de Economía.} \textbf{Elinor Ostrom} estudió cientos de casos de comunidades que administran un recurso común ---agua de riego, bosques, pesquerías--- y se propuso encontrar \textbf{cuál era el mejor conjunto de reglas}. \emph{No lo encontró: no hay reglas que sirvan en todas partes.} \textbf{Lo que sí encontró, y publicó en \emph{El gobierno de los bienes comunes} (1990), fueron \textbf{ocho principios de diseño} que aparecían en los sistemas que \emph{duraban} \textbf{y faltaban en los que se derrumbaban}.} \textbf{Tres de esos principios son exactamente el trabajo de hoy.} \emph{Uno:} \textbf{las reglas las hacen \emph{los propios usuarios}}, ajustadas a las necesidades y condiciones de su comunidad y su entorno. \emph{Dos:} \textbf{quienes vigilan el cumplimiento \emph{responden ante los usuarios}} del recurso ---no ante alguien de afuera---. \emph{Tres:} \textbf{tiene que haber mecanismos \emph{baratos} de resolución de conflictos} ---baratos en tiempo y en costo social: algo que la gente use de verdad---. \emph{Míralos juntos y verás lo mismo que hacen el CRIC y los misak, dicho por una economista que estudió medio mundo.}
\end{softbox}

\begin{softbox}{milcMostaza}{milcGris}{Pregunta-puente}
\textit{Ostrom encontró que \textbf{los sistemas que duran son aquellos cuyas reglas hacen los que las cumplen}. \textbf{¿Cuántas de las normas de tu salón las escribió alguien que se sienta ahí todos los días}\ldots \textbf{y cuántas de las que ustedes mismos propusieron han durado más de un mes}?}
\end{softbox}

\begin{softbox}{milcVerde}{milcGris}{Saber y saber hacer}
Al terminar podrás: (1) \textbf{explicar} por qué una norma hecha por quienes la van a cumplir tiene más posibilidades de sostenerse; (2) \textbf{redactar} reglas \textbf{verificables} y distinguirlas de los deseos; (3) \textbf{diseñar} quién verifica y ante quién se responde, con rendición pública; (4) \textbf{cerrar} el año en un acuerdo firmado que recoja los nueve momentos anteriores, con consecuencias y fecha de revisión.
\end{softbox}



\small
\begin{tabularx}{\textwidth}{>{\bfseries}p{3.0cm}Y}
\rowcolor{milcGradePrimary!12}Autor & Dr. Álvaro Cárdenas Orozco\\
DBA & Comprende los fenómenos que afectan a su territorio, regula sus emociones ante ellos y acompaña a otros con empatía (transversal 6°--11°, articulado con la política «Escuela, territorio de vida» del MEN, Resolución 006519 de 2025, y la Ley 1620 de 2013)\\
Referentes & Primera ayuda psicológica (OMS, 2011/2012) · Directrices IASC (2007) · USGS y Servicio Geológico Colombiano · Pueblo nasa y la minga (CRIC) · Dussel · Estoicismo · Floridi\\
Producto final & El acuerdo de aula del curso: entre cinco y ocho reglas escritas por ustedes a partir de los nueve momentos del año, cada una con lo que pasa si se rompe, quién lo verifica y cuándo se revisa —firmado por todos y con fecha de revisión puesta—.\\
\end{tabularx}
\normalsize

\puente{Este es el último momento del grado, y no trae tema nuevo: \textbf{convierte los nueve anteriores en reglas} que ustedes puedan usar el año entrante.}

\milcestacion{milcGradeDark}{Ruta MILC: así vamos a aprender}
\begin{tabularx}{\textwidth}{>{\centering\arraybackslash}X>{\centering\arraybackslash}X>{\centering\arraybackslash}X>{\centering\arraybackslash}X}
\phasecard{1}{milcVerde}{milcGris}{Escucha\\[-1mm]\small Observo, escucho y pregunto.} &
\phasecard{2}{milcTurquesa}{milcGris}{Sistematización\\[-1mm]\small Me doy la norma.} &
\phasecard{3}{milcMagenta}{milcGris}{Praxis\\[-1mm]\small Construyo y aplico.} &
\phasecard{4}{milcVino}{milcGris}{Evaluación\\[-1mm]\small Reflexión liberadora.}
\end{tabularx}


\puente{Primero, el inventario del año: qué aprendieron en cada momento que \emph{merece} volverse norma.}

\milcestacion{milcVerde}{Estación 1 de 3 · Escucha}

\begin{softbox}{milcVerde}{milcGris}{Escena para escuchar}
\textbf{Actividad 1 · IDENTIFICA --- Lo que el año dejó.} \textbf{Primera parte.} Repasa los nueve momentos del grado y escribe, para cada uno, \textbf{una cosa} que aprendiste y que \emph{le serviría a tu curso como norma}. \emph{Del 1: verificar antes de suponer. Del 2: los mandatos. Del 3 y el 4: el control y el límite. Del 5: las microagresiones. Del 6: las barreras. Del 7: el que llega. Del 8: la reparación colectiva. Del 9: reparar antes que castigar.} \textbf{Segunda parte.} De las normas que ya existen en tu salón ---las del manual, las del profesor, las que nadie escribió---, marca \textbf{cuáles se cumplen} y \textbf{cuáles no}, y al frente de las que no, escribe \textbf{por qué}. \emph{Casi siempre las razones son tres: nadie verifica, nadie sabe qué pasa si se rompe, o nadie estuvo de acuerdo con ella.} \textbf{Y la última:} de las cosas que su curso ha decidido hacer este año, \emph{¿cuántas siguen vivas?}
\end{softbox}

{\sffamily\bfseries\fontsize{8}{10}\selectfont\textcolor{milcNegro!55}{MARCA CUANDO LO HAYAS HECHO}}
\milccheckrow{Escribiste \textbf{una posible norma} salida de cada uno de los nueve momentos.}
\milccheckrow{Marcaste qué normas existentes se cumplen y cuáles no, con la razón.}
\milccheckrow{Contaste cuántas decisiones del curso siguen vivas.}

\begin{infoband}{milcVerde}


\textbf{En tu cuaderno (Actividad 1 · IDENTIFICA).} Título: «Actividad 1. Lo que el año dejó». \textbf{Formato:} las nueve posibles normas + las existentes con su razón de incumplimiento + cuántas decisiones siguen vivas. \textbf{Extensión:} media página. \textbf{Sabes que terminaste cuando} tienes \textbf{nueve candidatas} y sabes por qué fallan las normas que ya existen.
\end{infoband}

\puente{Ya lo tienen. Ahora, por qué las normas propias duran más, y qué hace que una regla se pueda comprobar.}

\milcestacion{milcTurquesa}{Fase 2 · Sistematización: entender la norma propia}

\begin{softbox}{milcTurquesa}{milcGris}{Cuatro claves sobre darse una norma}
Cuatro claves para que su acuerdo no termine como casi todos los acuerdos de curso: en una cartelera despegándose.
\end{softbox}

\milcpaso{1}{milcTurquesa}{\textbf{La hacen los que la cumplen.} Es el primer principio de Ostrom: \textbf{las reglas las elaboran \emph{los propios usuarios}}, ajustadas a sus necesidades y a sus condiciones (Ostrom, 1990). \emph{Y es exactamente lo que hacen los mandatos del CRIC y de los misak:} \textbf{la asamblea se da la norma}. \emph{La razón es práctica, no romántica:} \textbf{quien vive el problema sabe dónde está el detalle}. \emph{Una regla escrita desde afuera prohíbe lo que se ve desde afuera ---correr, gritar, el celular---, y casi nunca toca lo que de verdad hace daño en un salón:} \textbf{quién queda por fuera, qué se reenvía, cómo se arman los grupos.} \emph{Ustedes ya saben cuáles son ---llevan un grado entero estudiándolos---.}}
\milcpaso{2}{milcTurquesa}{\textbf{Pocas, claras y verificables.} \emph{Un acuerdo de veinte puntos no lo cumple nadie ---ni se acuerda---.} \textbf{Entre cinco y ocho es lo que funciona.} \emph{Y cada una tiene que pasar una prueba:} \textbf{¿cualquiera del curso podría decir, mirando, si se está cumpliendo?} \emph{Compara:} \textbf{«respetarnos» no se puede comprobar}; \emph{«los grupos de trabajo se arman por sorteo» sí; «nadie se sienta solo en el descanso más de dos días seguidos» sí; «lo que se dice en un chat privado no se trae al grupo» sí}. \textbf{La diferencia no es de tono: es que la segunda describe \emph{una conducta observable}.} \emph{Todo lo demás es un deseo, y los deseos no se incumplen ---simplemente no ocurren---.}}
\milcpaso{3}{milcTurquesa}{\textbf{Quién vigila, y ante quién responde.} Segundo principio de Ostrom: \textbf{quienes vigilan el cumplimiento \emph{rinden cuentas a los usuarios}}, no a alguien de afuera. \emph{Y es el par exacto de la rendición pública del CRIC:} \textbf{la autoridad responde ante la asamblea que la nombró}. \textbf{Aplicado a un curso, esto significa dos decisiones concretas:} \emph{quién verifica cada regla ---con nombre, y mejor si rota---} y \emph{ante quién informa}. \textbf{Y una advertencia:} \emph{si el único que verifica es el profesor, ustedes no se dieron una norma ---le pidieron una a alguien más---.} \textbf{Es la diferencia entre el bastón de mando y una orden.}}
\milcpaso{4}{milcTurquesa}{\textbf{Se revisa, o se muere ---y tiene que haber una salida barata---.} \emph{Tercer principio que usaremos:} \textbf{mecanismos \emph{baratos} de resolución de conflictos}, es decir, algo tan sencillo que la gente lo use de verdad en vez de dejar el problema crecer. \emph{En un curso puede ser tan simple como:} \textbf{«el que se sienta aludido lo dice en el momento y se para ahí»}, o \emph{«cinco minutos de revisión los viernes»}. \textbf{Y la revisión no es un adorno: es lo que hace que el acuerdo exista un mes después}, \emph{como aprendieron en el momento 10 del grado 8}. \textbf{La pregunta de cada revisión es siempre la misma:} \emph{¿se está cumpliendo? ¿quién no ha podido y por qué? ¿qué hay que ajustar?}}

\begin{drawbox}{milcTurquesa}{La forma de una regla que sirve}
\begin{pasoslista}
\item \textbf{Describe una conducta observable,} no una virtud. \emph{«Se arman por sorteo», no «somos incluyentes».}
\item \textbf{Dice qué pasa si se rompe} ---y no es un castigo: es qué se hace para arreglarlo---.
\item \textbf{Tiene quién verifica,} con nombre y rotando.
\item \textbf{Tiene fecha de revisión} ya puesta.
\item \textbf{Y el conjunto es corto:} \emph{entre cinco y ocho. Si son más, no las van a recordar.}
\end{pasoslista}
\end{drawbox}

\begin{softbox}{milcMostaza}{milcGris}{Errores comunes y su corrección}
\textbf{Reglas que son deseos:} «respetarnos», «ser unidos». \textbf{Veinte puntos:} nadie los recuerda. \textbf{Consecuencias que son castigos:} después de nueve momentos, ustedes saben que reparar funciona mejor. \textbf{Que solo verifique el profesor:} entonces no es un acuerdo propio. \textbf{Copiar el manual de convivencia:} ese ya existe; esto es lo que le falta. \textbf{No fechar la revisión:} sin fecha dura tres semanas. \textbf{Aprobar por mayoría sin oír a los que no están de acuerdo:} esos son los que la van a romper, y tienen razones.
\end{softbox}

\begin{infoband}{milcTurquesa}


\textbf{En tu cuaderno (Actividad 2 · EXPLICA).} Título: «Actividad 2. Deseo o regla». \textbf{Formato:} escribe \textbf{cinco} enunciados en su versión «deseo» y en su versión «regla verificable». \textbf{Extensión:} media página. \textbf{Sabes que terminaste cuando} puedes explicar por qué \textbf{un deseo no se incumple: simplemente no ocurre}.
\end{infoband}

\puente{Ya saben qué quieren. Es hora de escribirlo bien ---pocas, claras, con consecuencia y con responsable---.}

\milcestacion{milcMagenta}{Fase 3 · Praxis: escribir el acuerdo}

\begin{softbox}{milcMagenta}{milcGris}{Cómo se escribe un acuerdo que dure}
Este es el producto del año. Se escribe entre todos, se firma y se revisa: exactamente como un mandato.
\end{softbox}

\milcpaso{1}{milcMagenta}{\textbf{Lo que el año dejó (10 min, todo el curso).} Pongan en común las nueve candidatas y agrúpenlas. \emph{Van a ver que se repiten los mismos tres o cuatro asuntos ---quién queda por fuera, qué circula en los chats, cómo se arman los grupos, cómo se arregla lo que se rompe---.} \textbf{Escojan los cinco a ocho que de verdad les importan.} \emph{Y oigan con atención a quien esté en desacuerdo con alguno: es quien mejor sabe dónde va a fallar.}}
\milcpaso{2}{milcMagenta}{\textbf{Redactar las reglas (15 min).} Escríbanlas \textbf{una por una, en conducta observable}. \emph{Pasen cada una por la prueba:} \textbf{«¿cualquiera podría decir, mirando, si se cumple?»} \emph{Si la respuesta es no, reescríbanla ---no la descarten: casi siempre hay una versión verificable de lo que querían decir---.}}
\milcpaso{3}{milcMagenta}{\textbf{Consecuencia y verificación (12 min).} Para cada regla: \textbf{qué se hace si se rompe} ---y después del momento 9, ustedes ya saben que la consecuencia útil es \emph{reparadora} y no punitiva---; \textbf{quién verifica}, con nombre y \textbf{rotando}; y \textbf{ante quién se informa}. \emph{Añadan además el mecanismo barato:} \textbf{una frase que cualquiera pueda decir en el momento} para parar algo sin armar un problema.}
\milcpaso{4}{milcMagenta}{\textbf{Firmar y fechar (8 min).} \textbf{Fírmenlo todos} ---firmar no es un trámite: es lo que convierte una lista en un compromiso---, \textbf{péguenlo donde se vea} y \textbf{pongan la primera fecha de revisión ahora}, antes de salir del salón. \emph{Y decidan quién convoca esa revisión, con nombre.} \textbf{Después preséntenselo al director de grupo} ---no para que lo apruebe, sino para que lo conozca: \emph{esto es un mandato de ustedes, no un permiso que piden}---.}

\begin{drawbox}{milcMagenta}{Antes de cerrar el acuerdo}
\begin{checklista}
\item Son \textbf{entre cinco y ocho} reglas, no veinte.
\item Cada una describe una \textbf{conducta observable} y pasó la prueba de la verificación.
\item Cada una dice \textbf{qué se hace si se rompe}, y es reparador y no punitivo.
\item Cada una tiene \textbf{quién verifica}, con nombre y rotando, y ante quién informa.
\item Escucharon a quienes estaban en desacuerdo \textbf{antes} de aprobar.
\item Está \textbf{firmado}, pegado, y con \textbf{fecha de revisión y convocante} ya escritos.
\end{checklista}
\end{drawbox}

\begin{softbox}{milcMagenta}{milcGris}{Plantilla de trabajo}
\textbf{La regla.} \emph{«En este curso, \_\_\_\_.»} ---conducta observable---. \\[1mm]
\textbf{Lo demás.} \emph{«Si se rompe, \_\_\_\_. Verifica \_\_\_\_ (rota cada \_\_\_\_). Informa ante \_\_\_\_.»} \\[1mm]
\textbf{El cierre.} \emph{«Revisamos el \_\_\_\_. Convoca \_\_\_\_.»} ---y la frase barata: \emph{«el que se sienta aludido lo dice y se para ahí»}---.
\end{softbox}

\begin{infoband}{milcMagenta}


\textbf{En tu cuaderno (Actividad 3 · CREA).} Título: «Actividad 3. Nuestro acuerdo de aula». \textbf{Formato:} las cinco a ocho reglas completas ---conducta, consecuencia, verificador, instancia--- + la frase barata + la fecha de revisión y quién convoca. \textbf{Extensión:} una página. \textbf{Sabes que terminaste cuando} \textbf{está firmado y la fecha de revisión está escrita}.
\end{infoband}

\puente{El producto es un acuerdo firmado. \emph{Y su prueba no es la firma: es la primera revisión.}}

\milcestacion{milcMostaza}{Producto de la sesión}
\textbf{Acuerdo de aula: entre cinco y ocho reglas verificables, con consecuencia reparadora, verificador rotativo y fecha de revisión}.
\begin{softbox}{milcMostaza}{milcGris}{Criterios de calidad}
(1) El acuerdo tiene entre cinco y ocho reglas, redactadas como conductas observables. (2) Cada regla pasa la prueba de verificación: cualquiera podría decir si se cumple. (3) Las consecuencias son reparadoras y no punitivas, coherentes con el momento 9. (4) Hay verificador con nombre y rotación, y una instancia ante la cual se informa. (5) El acuerdo está firmado por todos y tiene fecha de revisión y convocante escritos.
\end{softbox}

\puente{Antes de cerrar, revisen la trampa clásica: si una regla no se puede comprobar, es un deseo con forma de norma.}

\milcestacion{milcVino}{Evaluación liberadora · cinco dimensiones}

\begin{softbox}{milcVino}{milcGris}{Sentido de la evaluación}
Evaluar no es solo revisar una nota. Es reconocer qué aprendí, qué cambió en mí, qué emociones manejé, cómo me relaciono con la ciudad y la comunidad, y qué le dejaré a quien viene detrás.
\end{softbox}

\textbf{Mi reflexión liberadora en cinco dimensiones.} Completa cada frase iniciada:

\small
\begin{tabularx}{\textwidth}{>{\bfseries\RaggedRight\arraybackslash}p{3.4cm}Y>{\RaggedRight\arraybackslash}p{4.6cm}}
\rowcolor{milcVino}\color{white}Dimensión & \color{white}\bfseries Mi respuesta (frame starter) & \color{white}\bfseries Algo que descubrí\\
1. Personal & Lo que más cambió en mí fue \linead{6.5cm} & \linead{4.2cm}\\
2. Emocional & La emoción más fuerte fue \linead{2.5cm} y la regulé \linead{3cm} & \linead{4.2cm}\\
3. Ciudadana & En democracia esto importa porque \linead{6cm} & \linead{4.2cm}\\
4. Local (Cartago) & En mi barrio sería útil para \linead{6.5cm} & \linead{4.2cm}\\
5. Intergeneracional & A mi abuela/abuelo le diría \linead{6.5cm} & \linead{4.2cm}\\
\end{tabularx}
\normalsize

\begin{infoband}{milcVino}
\textbf{Para la próxima sustentación.} Vuelve a esta tabla en seis meses. Verás que las respuestas cambian — no porque lo aprendido cambie, sino porque tú cambias. La evaluación liberadora es una conversación contigo a través del tiempo.
\end{infoband}


\puente{Cerramos el grado con tres voces: la comunidad que nace de escuchar, la coherencia entre decir y hacer, y el cuidado de lo compartido.}

\milcestacion{milcGradeDark}{Triángulo de pensamiento · cierre filosófico}

\begin{softbox}{milcVino}{milcGris}{Dussel · lente del nosotros}
\textit{«Aquel que es capaz de escuchar silenciosamente al Otro como otro es el que puede constituir una comunidad y no una mera sociedad totalizada.»}

{\footnotesize\itshape\color{milcNegro!70}--- Enrique Dussel · Filosofía de la liberación (1977)}

\textbf{Con esta frase se cierra el grado, y ahora se puede leer completa.} \emph{Una «sociedad totalizada» también puede tener reglas ---muy buenas reglas, incluso---:} \textbf{lo que no tiene es a alguien escuchando al que no piensa como la mayoría}. \emph{Por eso en la Actividad 3 hay una instrucción que parece menor y es la más importante:} \textbf{oír a quien esté en desacuerdo \emph{antes} de aprobar}. \textbf{No por cortesía ---por eficacia:} \emph{el que está en desacuerdo es el que sabe dónde va a fallar la regla, y además es el que la va a romper}. \emph{Un acuerdo aprobado por mayoría sobre las objeciones de una minoría no es un acuerdo: es una votación con la que hay que convivir.}

\textbf{Cuando decidimos algo en el curso, ¿alguien preguntó qué pensaban los que votaron en contra?}
\end{softbox}
\linead{\linewidth}\\[1mm]
\linead{\linewidth}

\begin{softbox}{milcMostaza}{milcGris}{Estoicismo · Séneca · Cartas a Lucilio, 20 (c. 64 d.C.) · cuidado interior}
\textit{«La filosofía enseña a obrar, no a hablar; quiere que cada uno viva a la manera que prescribe, que estén en armonía nuestras palabras con nuestras acciones.»}

{\footnotesize\itshape\color{milcNegro!70}--- Séneca · Cartas a Lucilio, 20 (c. 64 d.C.)}

\textbf{«Que estén en armonía nuestras palabras con nuestras acciones».} \emph{Eso es, exactamente, la diferencia entre un acuerdo y una cartelera.} \textbf{Y también es la razón por la que este acuerdo se firma:} \emph{la firma pone tu nombre al lado de una palabra que después vas a tener que sostener con hechos}. \emph{Séneca añade algo que este grado entero comprobó:} \textbf{la filosofía «enseña a obrar, no a hablar»}. \emph{Los nueve momentos anteriores no terminaron en opiniones ---terminaron en una conversación sostenida, un mandato desobedecido, una lista entregada, un protocolo, una reparación---.} \textbf{Este acuerdo es la última de esas obras, y la única que va a seguir funcionando cuando la clase termine.}

\textbf{De lo que firmé este año, ¿qué he sostenido con hechos?}
\end{softbox}
\linead{\linewidth}\\[1mm]
\linead{\linewidth}

\begin{softbox}{milcTurquesa}{milcGris}{Floridi · lente de la infoesfera}
\textit{«El florecimiento de las entidades informacionales y de toda la infoesfera debe promoverse preservando, cultivando y enriqueciendo sus propiedades.»}

{\footnotesize\itshape\color{milcNegro!70}--- Luciano Floridi · La ética de la información (2013; trad.)}

\textbf{Su acuerdo tiene que cubrir los dos lugares donde vive el curso: el salón y el chat.} \emph{Y el segundo casi siempre se olvida, aunque es donde ocurre la mitad de lo que trabajaron este año} ---lo que se reenvía, quién está y quién no, lo que queda escrito---. \textbf{Por eso conviene que al menos \emph{una} de sus reglas sea sobre el grupo:} \emph{qué no se trae ahí, qué se borra si alguien lo pide, quién revisa que estén todos los del curso}. \textbf{Y fíjate en el verbo de Floridi:} \emph{no dice «proteger» ---dice \textbf{cultivar} y \textbf{enriquecer}---}. \textbf{Un acuerdo que solo prohíbe cuida; uno que además dice qué sí se hace ahí, cultiva.} \emph{Los que duran son los segundos.}

\textbf{De nuestras reglas, ¿cuántas dicen qué sí hacemos, y cuántas solo qué no?}
\end{softbox}
\linead{\linewidth}\\[1mm]
\linead{\linewidth}

\begin{infoband}{milcNegro}
\textbf{Nota para el docente.} Las tres son citas textuales y llevan su referencia debajo. Puedes remitir a la obra en clase.
\end{infoband}

\milcestacion{milcVino}{Mi autoevaluación · marca una casilla por fila}

{\small
\setlength{\extrarowheight}{2.2pt}
\begin{tabularx}{\textwidth}{>{\RaggedRight\arraybackslash\bfseries}p{3.0cm}YYY}
\rowcolor{milcVino}\color{white}\bfseries Criterio & \color{white}\bfseries Lo logré & \color{white}\bfseries En proceso & \color{white}\bfseries Necesito apoyo\\[.6mm]
Comprensión del tema & \milcopcion{Explico las ideas con mis palabras.} & \milcopcion{Reconozco las ideas, falta precisión.} & \milcopcion{Necesito repasar lo básico.}\\[.8mm]
\rowcolor{milcGris}Calidad del acuerdo & \milcopcion{Nuestras reglas son pocas, verificables, con consecuencia, responsable y fecha de revisión.} & \milcopcion{Escribimos reglas, pero son deseos: no se puede comprobar si se cumplen.} & \milcopcion{Sigo pensando que las normas las pone el que manda y uno las cumple o no.}\\[.8mm]
Aplicación y producto & \milcopcion{Mi producto cumple los criterios.} & \milcopcion{Está armado, con ajustes pendientes.} & \milcopcion{Aún está en construcción.}\\[.8mm]
\rowcolor{milcGris}Reflexión 5 dimensiones & \milcopcion{Respondí cada una con honestidad.} & \milcopcion{Respondí algunas con detalle.} & \milcopcion{Mis respuestas son muy generales.}\\[.8mm]
Triángulo de pensamiento & \milcopcion{Conecté las 3 voces con mi trabajo.} & \milcopcion{Conecté al menos una.} & \milcopcion{No las relaciono con mi proceso.}\\[.8mm]
\end{tabularx}}

\vspace{3mm}
\textbf{Hoy entendí que:}\\[1mm]
\linead{\linewidth}\\[1mm]
\linead{\linewidth}

\textbf{Mi compromiso para la próxima sesión es:}\\[1mm]
\linead{\linewidth}\\[1mm]
\linead{\linewidth}

\begin{softbox}{milcMostaza}{milcGris}{Cita de despedida · Marco Aurelio}
\textit{``El obstáculo en el camino se vuelve el camino. Lo que estorba al camino se vuelve camino.''}\\[1mm]
Cada error de hoy, bien observado, se convierte en pieza del camino que pisarás mañana.
\end{softbox}

\small
\begin{tabularx}{\textwidth}{>{\bfseries}p{3.6cm}Y}
\rowcolor{milcGradeDark!12}Nombre completo: & \linead{10cm}\\
Fecha: & \linead{4cm} \quad Curso/grupo: \linead{4cm}\\
Mi firma: & \linead{9cm}\\
\end{tabularx}
\normalsize

\begin{infoband}{milcGradeDark}
\textbf{Para la primera revisión ---y para el año que viene.} \textbf{Primero, la revisión en la fecha que pusieron}, con las tres preguntas: \emph{¿se está cumpliendo? ¿quién no ha podido y por qué? ¿qué hay que ajustar?} \textbf{Y guarden el dato que importa:} \emph{cuántas de las reglas siguen vivas al mes}. \emph{Ese número, y no la firma, es la nota de este grado.} \textbf{Segundo, pregunta en tu casa por las normas propias:} averigua con alguien mayor \textbf{si perteneció alguna vez a algo que se diera sus propias reglas} ---una junta de acción comunal, un sindicato, una cooperativa, un grupo de la iglesia, un equipo, una asociación de padres---: \textbf{cómo las decidían, quién las hacía cumplir, qué pasaba con el que no las cumplía y si alguna vez tuvieron que cambiarlas}. \textbf{Trae:} el resultado de la revisión y lo que te contaron. \emph{Vas a encontrar los principios de Ostrom sin que nadie los haya leído: en los grupos que duraron, las reglas las hicieron los mismos socios, todos sabían quién vigilaba, y había una manera rápida y barata de arreglar los líos. En los que se acabaron, casi siempre falta alguna de las tres.}
\end{infoband}

\end{document}
